% TCC de teste no modelo do abnTeX2, escrito à mão, sem nada do AURA
% (passo 6.2.5, docs/latex-abntex.md §1.6).
\documentclass[12pt,openright,oneside,a4paper,chapter=TITLE,english,brazil]{abntex2}
\usepackage[utf8]{inputenc}
\usepackage{graphicx}
\usepackage{booktabs}

\titulo{Educação e diálogo na escola p\'ublica}
\autor{Maria da Silva \and João Souza}
\local{Vit\'oria}
\data{2025}
\orientador[Orientadora:]{Profa. Dra. Ana Lima}
\coorientador{Prof. Dr. Beto Reis}
\instituicao{Universidade Federal do Espírito Santo\par Centro de Educação\par Curso de Pedagogia}
\tipotrabalho{Trabalho de Conclusão de Curso}
\preambulo{Trabalho de Conclusão de Curso apresentado ao Curso de Pedagogia como requisito parcial.}

\begin{document}
\pretextual
\imprimircapa
\imprimirfolhaderosto

\begin{dedicatoria}
   \vspace*{\fill}
   \centering
   \noindent
   \textit{Aos meus pais.} \vspace*{\fill}
\end{dedicatoria}

\begin{agradecimentos}
Agradeço à orientadora.

E à banca.
\end{agradecimentos}

\begin{epigrafe}
    \vspace*{\fill}
	\begin{flushright}
		\textit{``Ensinar exige risco.''\\
		(Paulo Freire)}
	\end{flushright}
\end{epigrafe}

\begin{resumo}
 Este trabalho discute o diálogo na escola.

 \textbf{Palavras-chave}: educação. diálogo. escola.
\end{resumo}

\begin{resumo}[Abstract]
 \begin{otherlanguage*}{english}
   This work discusses dialogue at school.

   \vspace{\onelineskip}
   \noindent
   \textbf{Keywords}: education. dialogue.
 \end{otherlanguage*}
\end{resumo}

\begin{siglas}
  \item[ABNT] Associação Brasileira de Normas Técnicas
\end{siglas}

\tableofcontents*
\cleardoublepage

\textual

\chapter{Introdução}
\label{cap:intro}
A escola é espaço de diálogo \cite{freire}. Segundo Freire, ``ensinar exige risco'' \cite[p.~35]{freire}.

Os objetivos são:
\begin{itemize}
  \item compreender o diálogo;
  \item analisar a prática\footnote{Nota de rodapé.}.
\end{itemize}

\section{Metodologia}
Usamos $n = 30$ participantes, como mostra a Tabela~\ref{tab:dados}.

\begin{table}[htb]
\caption{Participantes por grupo}
\label{tab:dados}
\centering
\begin{tabular}{lc}
\toprule
Grupo & Total \\
\midrule
A & 12 \\
B & 18 \\
\bottomrule
\end{tabular}
\legend{Fonte: Dados da pesquisa.}
\end{table}

\begin{figure}[htb]
\caption{Fluxo da pesquisa}
\label{fig:fluxo}
\begin{center}
\includegraphics[width=0.8\textwidth]{imagens/fluxo}
\end{center}
\legend{Fonte: Elaborado pela autora.}
\end{figure}

\begin{citacao}
Citação longa de teste com mais de três linhas que a norma manda recuar \cite[p.~181]{freire}.
\end{citacao}

\begin{equation}
E = mc^2 \label{eq:einstein}
\end{equation}

\subsection{Análise}
Texto com \textsc{versalete} e acentos \'a \c{c}\~ao.

\chapter{Conclusão}
Conclusão do trabalho.

\postextual
\bibliography{referencias}

\begin{apendicesenv}
\partapendices
\chapter{Questionário}
Perguntas aplicadas.
\end{apendicesenv}

\begin{anexosenv}
\partanexos
\chapter{Lei de Diretrizes}
Texto da lei.
\end{anexosenv}

\printindex
\end{document}
